\section{Negative results about the spectrum of $\GPtwo$ sentences} \label{sec:gptwocounter}
The $1$-type spectrum of a $\Ctwo$ sentence is semilinear, and we showed how to calculate it in
the previous section. For $\GPtwo$, earlier in the paper we were able to characterize models by a linear system: but the solutions
to this linear system are not in one-to-one correspondence with the number of satisfiers of $1$-types.
Thus a natural question is whether the $1$-type spectrum of a $\GPtwo$ sentence is semilinear. We answer this below in
the negative.

Consider the following $\GPtwo$ formula $\varphi$ over vocabulary $\set{P_1, P_2, E}$:
\begin{equation*}
    \begin{aligned}
        \varphi\ :=\ &
        \forall x\ 
        \left(
            \left(\neg P_1(x) \land P_2(x)\right) \lor
            \left(P_1(x) \land \neg P_2(x)\right)
        \right)
        \ \ \land \ \ 
        \bigwedge_{1 \le i \le 2} \forall x\ P_i(x) \to \gamma_i(x) \\
        \gamma_1(x)\ :=\ &
        \presby{E(x, y) \land P_1(y)} - \presby{E(x, y) \land P_2(y)} = 0 \\
        \gamma_2(x)\ :=\ &
        \presby{E(y, x) \land P_1(y)} = 1.
    \end{aligned}
\end{equation*}
Let $\cG$ be a finite model of $\varphi$.
A vertex in $\cG$ is labeled with either $P_1$ or $P_2$.
Let $V_i$ be the set of vertices in $\cG$ with label $P_i$.
For every $v \in V_1$, the number of edges from $v$ to $V_1$
is the same as the number of edges from $v$ to $V_2$.
For every vertex $v \in V_2$, there exists exactly one edge from $V_1$ to $v$.
Therefore the cardinality of $V_2$ is the same as the number of edges between vertices of $V_1$,
which is bounded by $\abs{V_1}\cdot(\abs{V_1}-1)$.
Thus the 1-type spectrum of $\varphi$ is $\setc{(0, a, b, 0) \in \bbN^4}{b \le a(a-1)}$.
Hence we have the following negative result.
\begin{theorem}
    The 1-type spectrum of a $\GPtwo$ sentence is not always semilinear.
\end{theorem}

This does not resolve whether the set of cardinalities of finite models is a semilinear set, and
it leaves open whether $\GPtwo$ with global cardinality constraints is decidable.
